\begin{frame}
\frametitle{Kódování instrukcí procesoru RISC-V}

Má šest typů instrukcí:
\begin{table}
\footnotesize
\begin{tabular}{|m{0.4cm}|m{0.4cm}|m{1.0cm}|m{1.0cm}|m{0.4cm}|m{1.0cm}|m{1.0cm}|m{1.0cm}|m{0.4cm}|m{1.0cm}|}\hline
Typ & 31 & 30...25 & 24...21 & 20 & 19...15 & 14...12 & 11...8 & 7 & 6...0 \\ \hline
R & \multicolumn{2}{c|}{ fnct7 } & \multicolumn{2}{c|}{ rs2 } & rs1 & fnct3 &\multicolumn{2}{c|}{ rd } & opcode\\ \hline
I & \multicolumn{4}{c|}{ imm[11:0] } & rs1 & fnct3 &\multicolumn{2}{c|}{ rd } & opcode\\ \hline
S & \multicolumn{2}{c|}{ imm[11:5] } & \multicolumn{2}{c|}{ rs2 } & rs1 & fnct3 &\multicolumn{2}{c|}{ imm[4:0] } & opcode\\ \hline
B & imm [12] & imm [10:5]  & \multicolumn{2}{c|}{ rs2 } & rs1 & fnct3 & imm[4:1]& imm [11] & opcode\\ \hline
U & \multicolumn{6}{c|}{ imm[31:12] }  & \multicolumn{2}{c|}{ rd } & opcode\\ \hline
J & imm [20] & \multicolumn{2}{c|}{ imm[10:1] } & imm [11] & \multicolumn{2}{c|}{ imm[19:12] } & \multicolumn{2}{c|}{ rd } & opcode\\ \hline
\end{tabular}
\end{table}

\begin{itemize}
\item délka a typ instrukce (R,I,S,B,U,J) je určen 7-bity opcode
\item rs1,2 -- registery zdroje -- source; rd -- register cíle -- destination
\begin{itemize}
\item registry jsou vždy na pevné pozici v instrukci
\end{itemize}
\item immediate -- konstanty použité pro operaci nebo skoky
\item fnct3, fnct7 -- specifikace prováděné operace, např. sčítání, odčítání, posuny, atd.
\end{itemize}

\end{frame}

\begin{frame}
\frametitle{Kódování instrukcí procesoru RISC-V}

Význam hodnot opcode:
\begin{table}
\footnotesize
\begin{tabular}{|c|l|l|}\hline
Opcode & Skupina operací & Skutečná operace \\ \hline
0110011 & Typ R podle fnct & add, sub, slt, or, and \\ \hline
0010011 & Typ I podle fnct & addi, slt, or, and \\ \hline
0000011 & Čtení z paměti & lw \\ \hline
0100011 & Zápis do paměti & sw \\ \hline
1100011 & Větvení -- branch & beq \\ \hline
1101111 & Podprogramy & jal \\ \hline
1100111 & Návrat z podprogramu & jalr \\ \hline
0000111 & Načtení horní části konstanty & lui \\ \hline
\end{tabular}
\end{table}

Pro typ R jsou následující hodnoty \texttt{fnct3} a \texttt{fnct7}:
\begin{table}
\footnotesize
\begin{tabular}{|l|l|l|}\hline
fnct7 & fnct3 & Skutečná operace \\ \hline
0000000 & 000  & add \\ \hline
0100000 & 000  & sub \\ \hline
0000000 & 010  & slt \\ \hline
0000000 & 110  & bitový or \\ \hline
0000000 & 111  & bitový and \\ \hline
\end{tabular}
\end{table}

\end{frame}
